% 부록: 과제 2 그림 해설.
% 보고서(hw2/report)의 그림 아래 패널별 설명을 저자 지시(2026/09/30)로 이 부록에 옮겼다.
% 그림은 hw2/code/ 의 답별 스크립트가 만든 hw2/figures/*.png (씨드 20701313).
\chapter{부록: 과제 2 그림 해설}
\label{ch:hw2-figure-notes}

과제 2 보고서에는 그림과 짧은 캡션만 두고, 패널마다의 읽는 법은 이 부록에 모았다.
본문은 \Cref{ch:hw2}를 본다. 모든 모형 표본은 데이터와 같은 개수를 numpy 난수 생성기
(씨드 20701313)로 뽑았다. 보통 축(Linear Axes) 패널은 노선 0--200개만 보이고, 전체는
옆의 로그-로그 축(Log-Log Axes) 패널에 있다.

% =====================================================================
\section{공항 노선 수}
\label{sec:hw2app-airport}

\subsection{데이터의 분포}
\begin{figure}[H]
\centering
\includegraphics[width=0.95\linewidth]{p1a_0_data.png}
\caption{공항 노선 수의 경험적 분포($n = 3409$). 행은 PDF, PMF, CDF, CCDF이고 왼쪽은 보통 축,
오른쪽은 로그-로그 축이다. 초록 점선은 중앙값(4), 보라 점-선은 평균(19.8), 회색 파선은
$x = 60$이다.}
\label{fig:hw2app-airport-data}
\end{figure}
\begin{description}
  \item[PDF, 보통 축 (1행 왼쪽)] 막대 폭 5인 히스토그램(Histogram)으로 노선 0--200개를 보인다
    (200개 넘는 공항은 2.0\%). 공항의 56\%가 첫 막대(노선 5개 미만)에 있고, 82\%가 20개
    미만이며, 약 60개 너머의 막대는 바닥에 붙어 있다.
  \item[PDF, 로그-로그 (1행 오른쪽)] 로그 눈금에서 폭이 같은 구간으로 나누고 폭으로 나눈
    밀도(Density)다. 대략 직선으로 내려가다 약 100개부터 아래로 휜다.
  \item[PMF, 보통 축 (2행 왼쪽)] 노선 수 $k$마다 점 하나: 노선이 정확히 $k$개인 공항 수(왼쪽
    축)와 같은 높이를 비율 $P(X = k)$로 읽은 값(오른쪽 축). 노선 1개인 공항이 713곳(20.9\%),
    2개가 583곳, 3개가 363곳이다. 약 30개 이상은 모두 14곳 이하라 바닥에 깔린다.
  \item[PMF, 로그-로그 (2행 오른쪽)] 같은 점을 로그 축에 그렸다. 약 30개까지 기울기 약
    $-1.45$의 직선에 놓이며 $-\hat\alpha$에 가깝다. 파선($x = 60$) 오른쪽에서는 노선 수마다
    공항이 1, 2, 3곳뿐이라 점이 그 높이에 가로줄로 쌓인다. 분포의 모양이 아니라 개수의
    바닥이다.
  \item[CDF, 보통 축 (3행 왼쪽)] 거의 수직으로 올라 노선 60개에서 0.92, 200개에서 0.98에
    닿는다. 나머지 2\%는 축 밖에서 최댓값 915까지 이어진다.
  \item[CDF, 로그-로그 (3행 오른쪽)] $x = 1$에서 0.209로 뛰고 처음 몇 노선 수에서 계단으로
    오르다가 약 20개부터 1에 붙는다. 꼬리(Tail)가 거의 안 보여서 꼬리는 CCDF로 읽는다.
  \item[CCDF, 보통 축 (4행 왼쪽)] 노선 $x$개 이상인 공항의 비율이 1에서 노선 20개에서 0.18,
    60개에서 0.08로 떨어진다. 그 뒤 꼬리는 0 바로 위의 가는 선이다(200개에서 0.02).
  \item[CCDF, 로그-로그 (4행 오른쪽)] 약 100개까지 직선에 가깝다가 아래로 휘어 가장 큰 공항
    (노선 915개)에서 끝난다.
  \item[기준선] 평균이 중앙값의 5배 오른쪽에 있다. 노선이 수백 개인 허브 몇 곳이 평균을
    끌어올리기 때문이다. 로그-로그 CCDF에서 평균 이상인 공항은 약 18\%뿐이다.
  \item[요약] 중앙값 공항은 노선 4개이고 38\%가 1--2개뿐인데 가장 큰 공항은 915개다. 평균
    19.8은 중앙값의 5배이고 왜도(Skewness)는 6.0이다. 소수의 허브가 노선 대부분을 가진다.
\end{description}

\subsection{멱법칙 모형}
\begin{figure}[H]
\centering
\includegraphics[width=0.95\linewidth]{p1a_1_powerlaw.png}
\caption{공항 노선 수 대 멱법칙(Power Law, $\hat\alpha = 1.612$, $x_{\min} = 1$). 파랑은 데이터,
주황은 모형 표본이다. 행은 PDF, CDF, CCDF, QQ 플롯(QQ Plot), 왼쪽은 보통 축, 오른쪽은
로그-로그 축이다. $D = 0.209$.}
\label{fig:hw2app-airport-powerlaw}
\end{figure}
\begin{description}
  \item[PDF, 보통 축] 첫 5개 단위 막대들의 높이가 데이터와 모형에서 거의 같다. 200개 너머에
    데이터 2.0\%, 모형 표본 3.5\%가 있다.
  \item[PDF, 로그-로그] 약 300개까지 두 선이 같은 기울기의 직선이다. 그 뒤 데이터는 떨어져
    1000 근처에서 끝나고, 모형은 $10^5$ 너머까지 간다(가장 큰 표본 값은 342,563).
  \item[CDF, 보통 축] 200개까지 거의 겹친다. $x = 1$에서 데이터는 0.209(노선 1개인 공항
    20.9\%), 모형은 0.000이다. 연속 모형은 정확히 1에 확률을 두지 않는다.
  \item[CDF, 로그-로그] 로그 $x$축이 $x = 1, 2, 3$을 벌려 주어 가장 큰 간격이 $x = 1$(데이터
    0.209 대 모형 0)에 있음을 보여 준다. $D$ 자체는 보통 확률 단위로 $0.209 - 0 = 0.209$다.
    약 3개부터 두 곡선이 겹친다.
  \item[CCDF, 보통 축] 똑같아 보인다. 보통 축에서는 꼬리의 차이가 0에 눌린다.
  \item[CCDF, 로그-로그] 약 100개까지 같은 직선을 따른다. 그 뒤 데이터는 휘어 915에서 멈추고
    모형은 직선으로 계속 간다.
  \item[QQ, 보통 축] 축이 모형의 99.5\% 분위수(2,534)까지 늘어나 점 대부분이 왼쪽 아래 구석에
    몰린다. 모형의 거대 공항 몇 개만 눈에 띈다.
  \item[QQ, 로그-로그] 약 60개까지 $y = x$ 위에 있다가 선 아래 오른쪽으로 뻗는다. 99.5\%
    자리에서 모형은 2,534, 데이터는 358이다.
  \item[로그-로그 CCDF를 쓰는 이유] 멱법칙의 CCDF는 $P(X \ge x) = (x/x_{\min})^{-(\alpha-1)}$이고,
    양변에 로그를 씌우면 $\log P(X \ge x) = -0.612\,\log x$($\alpha = 1.612$, $x_{\min} = 1$)로
    기울기 $-(\alpha-1)$의 직선이 된다. CDF $1 - x^{-0.612}$는 ``$1 -$'' 때문에 로그가 풀리지
    않고, 로그-로그 축에서 모형과 데이터가 갈리는 꼬리에서 1에 붙는다.
\end{description}

\subsection{지수분포 모형}
\begin{figure}[H]
\centering
\includegraphics[width=0.95\linewidth]{p1a_2_exponential.png}
\caption{공항 노선 수 대 지수분포(Exponential, $\hat\lambda = 0.0504$). 배치는
\Cref{fig:hw2app-airport-powerlaw}와 같다. $D = 0.397$.}
\label{fig:hw2app-airport-exponential}
\end{figure}
\begin{description}
  \item[PDF, 보통 축] 모형의 첫 막대가 데이터보다 낮고, 노선 20--100개에 공항을 더 많이 둔다.
  \item[PDF, 로그-로그] 모형은 작은 공항에서 거의 평평하다가 약 50개부터 급히 떨어진다.
    데이터는 처음이 가파르고 꼬리가 길고 느리다. 표본의 6\%가 노선 1개 미만인데, 그런 공항은
    없다.
  \item[CDF, 보통 축] 모형은 약 190개에서 1에 닿는다(가장 큰 표본 값 191). 데이터는 200에서
    0.98이고 915까지 오른다.
  \item[CDF, 로그-로그] $x = 1$에서 데이터 0.209, 모형 0.056이다. 모형은 공항의 5.6\%를 노선
    1개 아래에 두고, 1에서 뛰지 않고 매끄럽게 오른다.
  \item[CCDF, 보통 축] 모형은 처음에 데이터보다 천천히 떨어지다가 약 190개에서 0이 된다.
    $x = 1$에서 데이터 1.000, 모형 0.944다.
  \item[CCDF, 로그-로그] 모형은 일찍 휘어 약 190개에서 끝나고, 데이터는 915까지 간다.
    $\lambda = 0.05$이면 노선 500개 공항의 확률은 약 $e^{-25}$인데, 데이터에는 그런 공항이
    5곳 있다.
  \item[QQ, 보통 축] 처음에만 $y = x$를 따르고 치솟는다. 99.5\% 자리에서 데이터는 358,
    모형은 약 100이다.
  \item[QQ, 로그-로그] 점이 너무 가파르다. 약 40개 아래에서는 $y = x$ 아래(작은 공항에 노선을
    너무 많이 줌), 그 위에서는 선 위(허브에 너무 적게 줌)다.
\end{description}

\subsection{균등분포 모형}
\begin{figure}[H]
\centering
\includegraphics[width=0.95\linewidth]{p1a_3_uniform.png}
\caption{공항 노선 수 대 균등분포(Uniform, $[\hat a, \hat b] = [1, 915]$). 배치는
\Cref{fig:hw2app-airport-powerlaw}와 같다. $D = 0.856$.}
\label{fig:hw2app-airport-uniform}
\end{figure}
\begin{description}
  \item[PDF, 보통 축] 모형은 높이 약 $1/914$의 평평한 선이고, 데이터의 첫 막대는 그 약 100배다.
    모형 표본의 78\%가 200개 너머에 있다(데이터는 2.0\%).
  \item[PDF, 로그-로그] 모형은 전 구간에서 $10^{-3}$ 근처로 평평하고, 데이터는 다섯 자릿수
    떨어진다.
  \item[CDF, 보통 축] 모형은 직선으로 올라 200개에서 약 0.20뿐이다. 데이터는 60개에서 0.92에
    닿는다. 가장 큰 간격 $D = 0.856$이 바로 노선 60개에 있다.
  \item[CDF, 로그-로그] $x = 1$에서 데이터 0.209, 모형 0이다. 모형은 약 90개까지 0.1 아래에
    있다.
  \item[CCDF, 보통 축] 모형은 직선으로 떨어져 200개에서도 약 0.80이고, 데이터는 왼쪽 끝에서
    무너진다.
  \item[CCDF, 로그-로그] 모형은 오른쪽 끝까지 1 근처에 있다가 915에서 수직으로 떨어진다.
  \item[QQ, 보통 축] 점이 바닥을 따라 평평하다. 모형 분위수가 약 900까지 오르는 동안 데이터는
    0--20에 머물다 맨 끝에서만 오른다.
  \item[QQ, 로그-로그] 모든 점이 $y = x$보다 한참 아래 오른쪽에 있다. 표본의 중앙값은 약
    470인데, 실제 중앙값은 4이고 실제 공항의 99\%가 285개 미만이다.
\end{description}

\subsection{정규분포 모형}
\begin{figure}[H]
\centering
\includegraphics[width=0.95\linewidth]{p1a_4_normal.png}
\caption{공항 노선 수 대 정규분포(Normal, $\hat\mu = 19.8$, $\hat\sigma = 53.5$). 배치는
\Cref{fig:hw2app-airport-powerlaw}와 같다. $D = 0.365$.}
\label{fig:hw2app-airport-normal}
\end{figure}
\begin{description}
  \item[PDF, 보통 축] 모형은 노선 20개 근처의 대칭 종 모양으로 음수 노선 수까지 퍼진다
    (0.1\% 분위수 약 $-136$). 데이터는 왼쪽 끝의 뾰족한 봉우리 하나다.
  \item[PDF, 로그-로그] 모형은 작은 공항에서 평평하다가 약 100개부터 급히 떨어진다. 표본의
    37\%가 노선 1개 미만(대부분 음수)이라 로그 축에 나오지 않는다.
  \item[CDF, 보통 축] 모형은 음수에서 시작해 $x = 1$에서 0.365(데이터 0.209)이고, 200개에서
    거의 1이다(가장 큰 표본 값 237). 데이터는 915까지 간다.
  \item[CDF, 로그-로그] 모형은 왼쪽 끝부터 약 10개까지 0.37 근처에서 평평하다. 그 평평한
    부분이 1 미만인 표본 37\%다.
  \item[CCDF, 보통 축] 모형은 S자로 0 근처에서 데이터와 교차한다. $x = 1$에서 모형 0.635,
    데이터 1.000이고, 200개에서 거의 0이다.
  \item[CCDF, 로그-로그] 약 100개 아래에서는 모형이 데이터 위에 있고, 꼬리는 약 230개에서
    끝나 915에 한참 못 미친다.
  \item[QQ, 보통 축] 모형의 낮은 분위수는 음수(약 $-118$까지)인데 데이터는 노선 1--3개다.
    약 100에서 $y = x$를 지나 치솟는다. 로그-로그 패널은 이 음수 분위수를 보여 주지 못한다.
  \item[QQ, 로그-로그] 지수분포처럼 너무 가파르다. 작은 공항에는 노선이 너무 많고 허브에는
    너무 적다. 정규분포는 $e^{-x^2}$로 줄어들어 꼬리가 얇다.
\end{description}

% =====================================================================
\section{영화 평점}
\label{sec:hw2app-movie}

\subsection{데이터의 분포}
\begin{figure}[H]
\centering
\includegraphics[width=0.95\linewidth]{p1b_0_data.png}
\caption{영화 평균 평점의 경험적 분포($n = 4392$). 위는 개수 히스토그램과 PDF, 아래는
CDF와 CCDF이며 모두 보통 축이다. 초록 점선은 중앙값(6.30), 보라 점-선은 평균(6.23)이다.}
\label{fig:hw2app-movie-data}
\end{figure}
\begin{description}
  \item[히스토그램 (1행 왼쪽)] 막대 폭 0.5의 영화 수. 6.25--6.75 구간에 약 1,000편으로 가장
    많고, 대부분(86\%)이 5--7.5점이다. 4점 미만은 1.3\%, 8점 이상은 1.5\%뿐이다.
  \item[PDF (1행 오른쪽)] 평점 값마다 막대 하나(평점은 0.1 단위로 반올림됨). 6.3 근처에 봉우리가
    하나이고 거의 대칭이다. 왼쪽 꼬리가 더 길다(왜도 $-0.48$). 2--4점 영화가 조금 있고
    8.5점을 넘는 영화는 없다.
  \item[PMF를 그리지 않은 이유] 평점은 여러 관객 점수의 평균이라 원래 연속인 양이고, 0.1
    단위로 반올림되어 기록되었을 뿐이다. PMF는 그 반올림 방식을 보여 줄 뿐이고(높이가
    반올림 단위에 따라 달라진다), 막대 폭 0.1인 PDF가 같은 정보를 밀도로 담는다. 참고로
    서로 다른 평점 값은 61개이고 가장 흔한 값은 6.5(209편)다.
  \item[CDF (2행 왼쪽)] 6.3 근처를 중심으로 한 S자 곡선이다(중앙값 6.3, 평균 6.23). 작은 계단은
    0.1 반올림 때문이다. 사분위수는 5.7과 6.8이다.
  \item[CCDF (2행 오른쪽)] CDF를 뒤집은 모양으로, ``이 평점 이상인 영화의 비율''을 읽는다.
    꼬리가 짧아서 로그 축이 필요 없다.
  \item[기준선] 평균이 중앙값보다 살짝 왼쪽에 있다. 아주 낮은 평점 몇 개가 평균을 끌어내리는,
    약간 왼쪽으로 치우친 표본에서 예상되는 모습이다.
\end{description}

\subsection{네 모형}
네 그림 모두 위는 PDF와 CDF(가장 큰 간격 $D$를 화살표로 표시), 아래는 QQ 플롯과
상자그림이다. 상자그림의 파선은 가능한 평점 1과 10이다.

\begin{figure}[H]
\centering
\includegraphics[width=0.95\linewidth]{p1b_1_powerlaw.png}
\caption{영화 평점 대 멱법칙($\hat\alpha = 1.851$, $x_{\min} = 1.9$). 왼쪽 PDF, 오른쪽 QQ
플롯(로그 축), 아래 오른쪽 상자그림(Box Plot). $D = 0.490$.}
\label{fig:hw2app-movie-powerlaw}
\end{figure}
\begin{description}
  \item[PDF] 모형은 가장 낮은 평점(1.9)에서 가장 높고, 데이터가 거의 없는 곳이다. 그 뒤로
    긴 꼬리가 이어지고, 표본의 14\%가 20을 넘는다. 가능한 평점이 아니다.
  \item[QQ] 표본의 99.5\% 분위수가 약 900이라 로그 축을 썼다. 높은 평점에서 점이 거의 수평으로
    오른쪽에 뻗는다. 모형의 높은 분위수는 수백인데 데이터는 8.5에서 멈춘다.
  \item[CDF] 가장 큰 간격 $D = 0.490$은 평점 4.68에 있다. 데이터는 4.7\%만 그 아래인데 모형
    표본은 53.8\%가 그 아래다.
  \item[상자그림] 표본의 24\%가 가능한 평점 1--10 밖(모두 10 초과)이고, 중앙값 4.2가 데이터의
    6.3보다 한참 낮다.
\end{description}

\begin{figure}[H]
\centering
\includegraphics[width=0.95\linewidth]{p1b_2_exponential.png}
\caption{영화 평점 대 지수분포($\hat\lambda = 0.1606$). $D = 0.481$.}
\label{fig:hw2app-movie-exponential}
\end{figure}
\begin{description}
  \item[PDF] 모형은 0에서 가장 높고 오른쪽 꼬리가 길다. 표본의 4\%가 20을 넘는다. 데이터는
    4--8.5 사이의 좁은 봉우리다.
  \item[QQ] 점이 $y = x$를 완만한 각도로 지난다. 모형의 퍼짐이 데이터보다 훨씬 크다. 낮은
    분위수는 0 근처, 높은 분위수는 20 너머다.
  \item[CDF] 가장 큰 간격 $D = 0.481$은 평점 4.69에 있다(데이터 4.7\%, 모형 52.8\%).
  \item[상자그림] 표본의 35\%가 1--10 밖이다(1 미만 15.9\%, 10 초과 19.3\%). 중앙값은 4.4다.
\end{description}

\begin{figure}[H]
\centering
\includegraphics[width=0.95\linewidth]{p1b_3_uniform.png}
\caption{영화 평점 대 균등분포($[1.9, 8.5]$). $D = 0.380$.}
\label{fig:hw2app-movie-uniform}
\end{figure}
\begin{description}
  \item[PDF] 모형은 $[1.9, 8.5]$에서 평평한데 데이터는 5--7.5에 몰려 있다.
  \item[QQ] 가운데에서 점이 평평하고, 왼쪽은 $y = x$ 위, 오른쪽은 아래다. 모형이 평점을 가운데
    쌓지 않고 고르게 펼쳐서, 낮은 분위수는 너무 낮고 높은 분위수는 너무 높다.
  \item[CDF] 모형은 직선으로 오르고 데이터는 S자다. 가장 큰 간격 $D = 0.380$은 평점 5.20에 있다.
  \item[상자그림] 표본은 1--10 안에 있지만, 상자(사분위수 3.6--6.9)가 데이터(5.7--6.8)의 두 배
    넘게 넓다.
\end{description}

\begin{figure}[H]
\centering
\includegraphics[width=0.95\linewidth]{p1b_4_normal.png}
\caption{영화 평점 대 정규분포($\hat\mu = 6.227$, $\hat\sigma = 0.893$). $D = 0.048$.}
\label{fig:hw2app-movie-normal}
\end{figure}
\begin{description}
  \item[PDF] 표본이 데이터와 같은 종 모양을 그리고 봉우리 위치도 같다.
  \item[CDF] 두 곡선이 거의 겹친다. 가장 큰 간격 $D = 0.048$은 평점 6.30에 있다.
  \item[상자그림] 표본의 상자와 수염이 데이터와 거의 같다.
  \item[QQ] 약 4.5--7.5에서 $y = x$ 위에 있다. 양 끝에서만 벗어난다. 데이터의 가장 낮은
    평점이 모형 예측보다 낮고(긴 왼쪽 꼬리), 데이터는 8.5에서 멈추는데 모형은 계속된다.
    $D = 0.048$로 가장 작다.
\end{description}

% =====================================================================
\section{배우 나이}
\label{sec:hw2app-age}

\begin{figure}[H]
\centering
\includegraphics[width=0.8\linewidth]{p2_0_age_all.png}
\caption{1--99세 필터 전의 모든 나이(2026년 기준, $n = 460{,}589$). 파랑은 남기는 나이,
코랄은 지우는 나이이며 세로축은 로그 눈금이다. 초록 점선은 중앙값(Median), 보라 점쇄선은 평균(Mean)이다.}
\label{fig:hw2app-age-all}
\end{figure}
\begin{description}
  \item[남김 (파랑)] 27--99세, 453,089명. 과제가 1--99세를 남기라고 한다. 가장 늦은 출생 연도가 1999년이라 27세 미만은 없다.
  \item[지움 (코랄)] 100--294세, 7,500명(1.6\%). 대부분 100세에서 약 130세이고, 150세,
    165세, 294세(1732년생)에 외딴 막대가 있다. 2015년 사진 데이터베이스의 사람으로 가능한
    나이가 아니고 과제 범위 밖이라 필터가 지운다. 연도가 \texttt{0000}인 129행은 나이가 아니라서 이 그림
    전에 빠진다.
  \item[중앙값과 평균] 필터 전 전체: 중앙값 55세, 평균 57.65세. 지운 100세 이상의 꼬리가
    평균만 끌어올린다.
  \item[로그 세로축인 이유] 지우는 쪽이 1.6\%뿐이라, 보통 축이면 연도당 약 16,000명인
    봉우리 옆에서 막대가 보이지 않는다.
\end{description}

\begin{figure}[H]
\centering
\includegraphics[width=0.8\linewidth]{p2_1_age.png}
\caption{2026년 기준 나이($n = 453{,}089$). 초록 점선은 중앙값(Median), 보라 점쇄선은 평균(Mean)이다.}
\label{fig:hw2app-age}
\end{figure}
\begin{description}
  \item[모양] 넓은 봉우리 하나: 27세에서 시작해 30대에 오르고 40--50대에 봉우리, 99세까지
    긴 오른쪽 꼬리. 대부분 36--71세 사이다.
  \item[들쭉날쭉한 막대] 이웃한 나이끼리 높이가 다르다(예: 48세와 57세가 두드러짐). 해마다
    태어난 배우 수의 잡음이지 나이 자체의 규칙이 아니다.
  \item[중앙값과 평균] 중앙값 55세, 평균 56.79세. 평균이 중앙값보다 크므로 오른쪽 꼬리가 길다.
\end{description}

\begin{figure}[H]
\centering
\includegraphics[width=0.8\linewidth]{p2_2_first_digit.png}
\caption{나이의 첫째 자리 숫자. 회색 파선은 균등 수준 $1/9$,
초록 파선은 자릿수에 맞춘 정규분포($\mu = 5.22$, $\sigma = 1.43$)다. 초록 점선은 중앙값(Median), 보라 점쇄선은 평균(Mean), 금색 실선은 최빈값(Mode)이고, 옅은 보라 띠는 평균 $\pm$ 표준편차(Standard Deviation) 구간이다.}
\label{fig:hw2app-first-digit}
\end{figure}
\begin{description}
  \item[막대] 숫자 1--9의 비율은 0\%, 0.7\%, 7.9\%, 24.5\%, 30.0\%, 18.9\%, 10.0\%, 5.4\%,
    2.5\%다. 10--19세가 없으므로 숫자 1은 하나도 없다. 나이 히스토그램의 봉우리와 같은 5와 4에서 가장 높다.
  \item[균등선] $1/9 = 0.111$과 멀다. 4--6은 한참 위, 1--3과 8--9는 한참 아래다.
  \item[정규 곡선] 자릿수의 평균과 표준편차로 맞춘 정규분포를 숫자 $d$마다 $d \pm 0.5$ 넓이로
    그렸다. 중심과 폭은 대략 맞지만 최대 5.3\%포인트 어긋난다(숫자 4). 데이터는 4와 5에서 더 뾰족하고
    8과 9가 더 많다(왜도 $+0.53$). 첫째 자리는 나이 분포를 10년 단위로 묶은 것이고 나이가
    오른쪽으로 치우쳐 있기 때문이다.
  \item[중앙값과 평균] 중앙값 5, 평균 5.22. 평균이 조금 오른쪽에 있다(왜도 $+0.53$과 같은 방향).
  \item[최빈값] 5(30.0\%). 중앙값 5와 같고 평균 5.22는 조금 오른쪽이다. 셋이 거의 모이므로
    봉우리 하나인 분포이고, 평균만 오른쪽 꼬리 쪽으로 끌려 있다.
  \item[분산과 표준편차] 분산(Variance) 2.05, 표준편차 $\sigma = 1.43$. 평균 $\pm\,\sigma$ 구간은
    3.79--6.65로, 숫자 4, 5, 6이 그 안에 든다.
\end{description}

\begin{figure}[H]
\centering
\includegraphics[width=0.8\linewidth]{p2_3_last_digit.png}
\caption{나이의 마지막 자리 숫자. 파선은 균등 수준 $1/10$이다. 초록 점선은 중앙값(Median), 보라 점쇄선은 평균(Mean), 금색 실선은 최빈값(Mode)이고, 옅은 보라 띠는 평균 $\pm$ 표준편차(Standard Deviation) 구간이다.}
\label{fig:hw2app-last-digit}
\end{figure}
\begin{description}
  \item[막대] 모든 숫자가 8.97\%(숫자 3)--11.12\%(숫자 7) 사이로, 10\%에서 1.12\%포인트
    안이다.
  \item[균등선] 열 막대 모두 0.100 파선 근처에 있고, 두드러지게 많은 숫자가 없다.
  \item[중앙값과 평균] 중앙값 5, 평균 4.55. 완전 균등한 0--9라면 평균은 4.5이므로 거의 같다.
  \item[최빈값] 7(11.12\%). 다른 숫자보다 1\%포인트쯤 많을 뿐이다. 균등분포에서는 모든 숫자가
    최빈값이므로 여기서 최빈값은 잡음이 고른 숫자다. 평균 4.55, 중앙값 5, 최빈값 7이 서로
    다른 것은 봉우리가 없다는 뜻이다.
  \item[분산과 표준편차] 분산 8.24, 표준편차 $\sigma = 2.87$. 완전 균등한 0--9의 분산은
    $(10^2 - 1)/12 = 8.25$이므로 이것도 거의 같다. 첫째 자리($\sigma = 1.43$)보다 두 배 퍼져 있다.
\end{description}
